% Part: methods
% Chapter: induction
% Section: inductive-definitions

\documentclass[../../../include/open-logic-section]{subfiles}

\begin{document}

\olfileid{mth}{ind}{idf}

\olsection{আরোহী সংজ্ঞা}

যুক্তিবিদ্যায় আমরা প্রায়ই কোনো শ্রেণির বস্তুকে \emph{আরোহীভাবে}
সংজ্ঞায়িত করি। অর্থাৎ, কোনগুলি সেই শ্রেণির বস্তু হবে তা
বলে দেওয়ার পাশাপাশি ওই শ্রেণির পুরোনো বস্তু থেকে নতুন
বস্তু পাওয়ার নিয়মও নির্দিষ্ট করি। যেমন, কোনো ভাষার পদ ও
!!{formula}s-এর মতো বিশেষ প্রতীক-অনুক্রম প্রায়ই এভাবে
সংজ্ঞায়িত করা হয়। সহজ উদাহরণ হিসেবে $\mathrm{a}$,
$\mathrm{b}$, $\mathrm{c}$, $\mathrm{d}$ অক্ষর,~$\circ$ প্রতীক,
এবং $[$ ও $]$ বন্ধনী দিয়ে গঠিত স্ট্রিংগুলি ভাবো। যেমন,
``$[[\mathrm{c} \circ \mathrm{d}][$'', ``$[\mathrm{a}[]\circ]$'',
``$\mathrm{a}$'' অথবা ``$[[\mathrm{a} \circ \mathrm{b}]\circ
\mathrm{d}]$''। প্রথম দুটি স্ট্রিংয়ে নিশ্চয় কিছু ``গোলমাল''
লাগছে: প্রথমটিতে বন্ধনীগুলি মেলেনি; দ্বিতীয়টিতে মনে হয়
``$\circ$'' এমন দুটি অভিব্যক্তিকে ``জোড়া'' উচিত, যেগুলির
নিজেদেরও অর্থপূর্ণ গঠন আছে। তৃতীয় ও চতুর্থ স্ট্রিংটি ভালো
লাগে: প্রত্যেক ``$[$''-এর জন্য একটি মিলন্ত ``$]$'' আছে
(বন্ধনী থাকলে), আর প্রত্যেক $\circ$-এর দুই পাশে জোড়া
বন্ধনী-ঘেরা ``সঠিক'' অভিব্যক্তি পাওয়া যায়।

কোনটিকে ``সুগঠিত পদ'' বলব, তা নির্ভুলভাবে স্থির করতে চাই।
প্রথমত, প্রত্যেকটি অক্ষর একাই সুগঠিত। শুধু একটি অক্ষর নয়
এমন যেকোনো সুগঠিত পদ ``$[t \circ s]$'' আকারের হবে,
যেখানে $s$ ও $t$ নিজেরাও সুগঠিত। উল্টো দিকে, $t$ ও $s$
সুগঠিত হলে তাদের মাঝে $\circ$ বসিয়ে এবং দুদিকে একজোড়া
বন্ধনী দিয়ে নতুন সুগঠিত পদ গড়া যায়। এই ক্রিয়াগুলির
সাহায্যেই সুগঠিত পদগুলির সেট \emph{সংজ্ঞায়িত} করতে পারি।
এটি একটি \emph{আরোহী সংজ্ঞা}।

\begin{defn}[সুগঠিত পদ]
  \emph{সুগঠিত পদ}-এর সেট নিচের নিয়মে আরোহীভাবে সংজ্ঞায়িত:
  \begin{enumerate}
  \item $\mathrm{a}$, $\mathrm{b}$, $\mathrm{c}$,
    $\mathrm{d}$—প্রতিটি অক্ষর একটি সুগঠিত পদ।
  \item $s_1$ ও $s_2$ সুগঠিত পদ হলে
    $[s_1 \circ s_2]$-ও সুগঠিত পদ।
  \item আর কিছু সুগঠিত পদ নয়।
  \end{enumerate}
\end{defn}

এই সংজ্ঞা বলে, কোনো বস্তু সুগঠিত পদ হবে যদি এবং কেবল
যদি সসীমসংখ্যক ধাপে (1) ও~(2) নিয়ম মেনে সেটি গড়া যায়।
প্রথম ধাপে কেবল একেকটি অক্ষর নিয়ে সুগঠিত পদগুলি গড়ি:
\[
\mathrm{a}, \mathrm{b}, \mathrm{c}, \mathrm{d}
\]
দ্বিতীয় ধাপে ইতিমধ্যে গড়া পদগুলির উপর (2) প্রয়োগ করি। পাই:
\[
[\mathrm{a} \circ \mathrm{a}], [\mathrm{a} \circ \mathrm{b}],
[\mathrm{b} \circ \mathrm{a}], \dots, [\mathrm{d} \circ \mathrm{d}]
\]
অর্থাৎ দুটি অক্ষরের সম্ভাব্য সব জোড়া। তৃতীয় ধাপে আগের
যেকোনো দুটি সুগঠিত পদের উপর আবার (2) প্রয়োগ করি।
পাই যেমন $[\mathrm{a} \circ [\mathrm{a} \circ
    \mathrm{a}]]$—এখানে $t$ হলো প্রথম ধাপের $\mathrm{a}$
এবং $s$ হলো দ্বিতীয় ধাপের $[\mathrm{a} \circ \mathrm{a}]$।
আরও পাই $[[\mathrm{b} \circ
    \mathrm{c}] \circ [\mathrm{d} \circ \mathrm{b}]]$,
যা দ্বিতীয় ধাপের $[\mathrm{b} \circ \mathrm{c}]$ এবং $[\mathrm{d}
  \circ \mathrm{b}]$ দিয়ে গড়া। এভাবেই এগোয়। (3) নম্বর
দফাটি এই নিয়মে গড়া হয়নি এমন কোনো বস্তু সুগঠিত পদের
সেটে ঢুকতে দেয় না।

লক্ষ করো, দেখতে ``সঠিক'' লাগে এমন প্রত্যেকটি প্রতীক-অনুক্রম
এই সংজ্ঞা অনুযায়ী সুগঠিত—এটি এখনো প্রমাণ করিনি। তবে
সংজ্ঞা মেনে যা গড়তে পারি তা যে দেখতে ``সঠিক'' লাগে, তা
স্পষ্ট: বন্ধনীগুলি মেলে এবং $\circ$ নিজ নিজ গঠনে সঠিক
দুটি অংশকে যুক্ত করে।

আরোহী সংজ্ঞার প্রধান সুবিধা হলো, সব সুগঠিত পদ সম্পর্কে
কিছু প্রমাণ করতে চাইলে সংজ্ঞাই বলে দেয় কোন ক্ষেত্রগুলি
বিবেচনা করতে হবে। যেমন, $t$ একটি সুগঠিত পদ জানলে
আরোহী সংজ্ঞা থেকে $t$-এর সম্ভাব্য গঠন জানা যায়: $t$
একটি অক্ষর হতে পারে, অথবা কোনো সুগঠিত পদ $s_1$ ও~$s_2$
দিয়ে $[s_1 \circ s_2]$ হতে পারে। (3) নম্বর দফার কারণে
এই দুটি ছাড়া আর কোনো সম্ভাবনা নেই।

আরোহীভাবে সংজ্ঞায়িত সেটের সব বস্তুর বিষয়ে প্রমাণে
প্রবল আরোহ বিশেষভাবে গুরুত্বপূর্ণ। ধরা যাক, প্রমাণ করতে
চাই যে দৈর্ঘ্য~$n$-এর প্রত্যেক সুগঠিত পদে $[$-এর সংখ্যা
$< n/2$। এটিকে সব~$n$-এর দাবি হিসেবেও দেখা যায়:
প্রত্যেক $n$-এর জন্য দৈর্ঘ্য~$n$-এর যেকোনো সুগঠিত পদে
$[$-এর সংখ্যা $< n/2$।

\begin{prop}
  যেকোনো $n$-এর জন্য, দৈর্ঘ্য~$n$-এর সুগঠিত পদে
  $[$-এর সংখ্যা $< n/2$।
\end{prop}

\begin{proof}
এই ফল (প্রবল) আরোহে প্রমাণ করতে নিচের শর্তাধীন দাবিটি
দেখাতে হবে:
\begin{quote}
  প্রত্যেক $l < k$-এর জন্য দৈর্ঘ্য~$l$-এর যেকোনো সুগঠিত
  পদে $[$-এর সংখ্যা $< l/2$ হলে, দৈর্ঘ্য~$k$-এর যেকোনো
  সুগঠিত পদে $[$-এর সংখ্যা $< k/2$।
\end{quote}
শর্তাধীন বক্তব্যটি দেখাতে তার পূর্বপক্ষ সত্য ধরি: যেকোনো
$l<k$-এর জন্য দৈর্ঘ্য~$l$-এর সুগঠিত পদে $[$-এর সংখ্যা
$< l/2$। এটিই আরোহের অনুমান। দেখাতে চাই, দৈর্ঘ্য~$k$-এর
সুগঠিত পদের ক্ষেত্রেও একই দাবি সত্য।

ধরা যাক, $t$ হলো দৈর্ঘ্য~$k$-এর একটি সুগঠিত পদ। পদগুলি
আরোহীভাবে সংজ্ঞায়িত বলে দুটি ক্ষেত্র: (1)~$t$ নিজে একটি
অক্ষর, অথবা (2)~কোনো সুগঠিত পদ $s_1$ ও~$s_2$ দিয়ে
$t$ হলো $[s_1 \circ s_2]$।
\begin{enumerate}
\item $t$ একটি অক্ষর। তাহলে $k = 1$ এবং $t$-তে
$[$-এর সংখ্যা~$0$। যেহেতু $0 < 1/2$, দাবি সত্য।
\item কোনো সুগঠিত পদ $s_1$ ও~$s_2$ দিয়ে $t$ হলো
  $[s_1 \circ s_2]$। ধরি, $s_1$-এর দৈর্ঘ্য $l_1$ এবং
  $s_2$-এর দৈর্ঘ্য $l_2$। তাহলে $t$-এর দৈর্ঘ্য~$k$
  হলো $l_1+l_2+3$ ($s_1$, $s_2$-এর দৈর্ঘ্যের সঙ্গে
  $[$, $\circ$, $]$—এই তিনটি প্রতীক)। $l_1+l_2+3$
  সব সময় $l_1$-এর চেয়ে বড়, তাই $l_1 < k$। একইভাবে
  $l_2 < k$। ফলে আরোহের অনুমান $s_1$ ও~$s_2$-এর
  জন্য প্রযোজ্য: $s_1$-এ $[$-এর সংখ্যা~$m_1$ হলো
  $< l_1/2$, আর~$s_2$-এ $[$-এর সংখ্যা~$m_2$ হলো
  $< l_2/2$।

  $t$-তে $[$-এর সংখ্যা হলো~$s_1$-এ $[$-এর সংখ্যা,
  $s_2$-এ $[$-এর সংখ্যা এবং আরও~$1$ যোগ করে—অর্থাৎ
  $m_1 + m_2 + 1$। যেহেতু $m_1 < l_1/2$ ও $m_2 < l_2/2$,
  তাই
  \[
  m_1 + m_2 + 1 < \frac{l_1}{2} + \frac{l_2}{2} + 1 = \frac{l_1+l_2+2}{2} < \frac{l_1+l_2+3}{2} = k/2.
  \]
\end{enumerate}
দুই ক্ষেত্রেই (আরোহের অনুমান ব্যবহার করে) দেখিয়েছি,
$t$-তে $[$-এর সংখ্যা $< k/2$। প্রবল আরোহে প্রস্তাবটি
অনুসৃত।
\end{proof}

\begin{prob}
  অতিসুগঠিত পদগুলির সেট নিচের নিয়মে সংজ্ঞায়িত করো:
  \begin{enumerate}
  \item $\mathrm{a}$, $\mathrm{b}$, $\mathrm{c}$,
    $\mathrm{d}$—প্রতিটি অক্ষর একটি অতিসুগঠিত পদ।
  \item $s$ অতিসুগঠিত পদ হলে $[s]$-ও তাই।
  \item $s_1$ ও $s_2$ অতিসুগঠিত পদ হলে
    $[s_1 \circ s_2]$-ও অতিসুগঠিত পদ।
  \item আর কিছু অতিসুগঠিত পদ নয়।
  \end{enumerate}
  দেখাও যে দৈর্ঘ্য~$n$-এর অতিসুগঠিত পদ~$t$-তে
  $[$-এর সংখ্যা $\le n/2 +1$।
\end{prob}

\end{document}
